% Sequences: reduce is free to group the combines any way, which is what makes
% the span logarithmic. Part (i) makes them exhibit a grouping rather than just
% assert "it's parallel"; (ii) then reads the bound off their own answer.
% Associativity is only a closing remark here -- the question does not ask for
% it, so the refsol should not grade on it.
%
% Verified in SML/NJ: balanced grouping of op+ over 1..n has combine-tree depth
% exactly ceil(log2 n) -- 2 for n=4, 3 for n=8, 10 for n=1000.
% Non-associativity check: (op -) gives -10 left-nested vs 0 balanced.
\part[4]\hfill

  We are interested in the evaluation of the expression
  \code{Seq.reduce (op+) 0 S}.

  \begin{enumerate}
    \item Write an expression of \code{+} operations which is extensionally equivalent
    to the above, and represents a realistic way that it would actually be evaluated.

    \item Explain why this expression gives an advantageous performance characteristic
    to the original expression.
  \end{enumerate}

  \begin{solutionorbox}[19em]
    \textbf{(i)} A \textbf{balanced} grouping: split the sequence in half,
    reduce each half, and add the two results. On
    \code{S} $= \langle s_0, s_1, s_2, s_3 \rangle$ that is
    $$(s_0 + s_1) + (s_2 + s_3)$$
    rather than the left-nested $(((0 + s_0) + s_1) + s_2) + s_3$. Any answer
    whose parenthesization is balanced is fine --- the point is that the two
    halves appear as independent subexpressions.

    \textbf{(ii)} The two halves are independent, so they can be evaluated in
    \textbf{parallel}: the span is one $+$ plus the span of the deeper half,
    $$S(n) = S(n/2) + O(1) = O(\log n)$$
    A balanced grouping over $n$ elements is a combine tree of depth
    $\lceil \log_2 n \rceil$ --- each level halves the number of partial results
    ($n \to n/2 \to n/4 \to \cdots \to 1$) --- so only $\log n$ rounds are
    needed. The \textbf{work} is unchanged, since all $n-1$ additions still
    happen; what changes is that they no longer have to happen one after
    another. The left-nested grouping makes each addition depend on the previous
    one, a chain of depth $n$, giving span $O(n)$.

    \vspace{4pt}
    (This is also why \code{reduce} requires an \textbf{associative} combiner:
    it may pick any grouping, so all groupings must agree. \code{op+} qualifies;
    \code{op-} would not.)
  \end{solutionorbox}
